% page 129 du fac-simile (image p-128 du scan)
% bloc 165.1 x 250.6 mm ; marge gauche 36.9 mm
\pgc{15.240mm}{0.000mm}{
\l{}
\l{}
\l{\cel{54}119}
\l{}
\l{}
\l{dominanto. I. (\sou{muziko}). En la gamo ordinara, la noto kinesma,}
\l{\cel{3}ta qua (pro formacar kun la unesma la konsonanco maxim per-}
\l{\cel{3}fekta) determinas la tono ed uzesas kom fundamento di la har-}
\l{\cel{3}monio.\cel{2}II. En la plankanto, noto sur qua onu kantas la korpo}
\l{\cel{3}di verso.\cel{2}III. (\sou{kristalografi}o) Tipo geometriala normala}
\l{\cel{3}segun qua korpo kristaleskas.- DEF.}
\l{}
\l{\sou{dominikano}.\cel{2}Reguliero de la ordeno di la frati predikera, fon-}
\l{\cel{3}dita da Santa Dominikus. - DEFIRS.}
\l{}
\l{\sou{domtar}. (\sou{trans.})\cel{2}I. Koaktar ad obedio (animalo ne-amansita.}
\l{\cel{3}- II. Koaktar ad obedio (ti qui refuzas submisar su). - DFIS.}
\l{}
\l{\sou{donar}. (\sou{trans., ad)}\cel{2}I. Igar disponebla da ulu. - II. Propozar}
\l{\cel{3}kondiciono determinita en la enunco di problemo. - EFIS.}
\l{}
\l{\sou{donacar}. (\sou{trans}.)Abandonar ad ulu la proprieto di ulo, sen re-}
\l{\cel{3}cevar ulo po olu. - EFIRS.}
\l{}
\l{\sou{dop}.\cel{2}An la latero opozita ad olta ube situesas la vizajo di}
\l{\cel{3}persono o la facio di kozo. - I.}
\l{}
\l{\sou{dorado}.\cel{2}(\sou{zool.}) Mar-fisho de la familio \textquotedbl{}komberidi\textquotedbl{}. - L.\cel{1}\sou{Co-}}
\l{\cel{3}\sou{ryphaena hippurus}.- II. Mar-fisho de la familio \textquotedbl{}sparoidi\textquotedbl{}.}
\l{\cel{3}L.\cel{1}\sou{Chrysophrys aurata}. - DFIRS.}
\l{}
\l{\sou{dorika}. (\sou{arkitekturo}). La duesma ek la kin grupi arkitekturala}
\l{\cel{3}che la Antiqui, karakterizata da koloni kaneloza, sen-baza,}
\l{\cel{3}kapitelo ek granda muluro kalicatra, fuzo interruptata da}
\l{\cel{3}triglifi. - DEFIRS.}
\l{}
\l{\sou{dorlotar}. (\sou{ulu)}. Traktar tenere, afablege. - F.}
\l{}
\l{\sou{dormar}. (\sou{netrans.})\cel{2}Esar en ta stando dum qua (che homo od ani-}
\l{\cel{3}malo, haltigesas ula funcioni di la vivo (precipue dum nokto),}
\l{\cel{3}pro la neceseso di lia repozo. - EFIS.}
\l{}
\l{\sou{dorno}.\cel{1}\sou{(bot}.)Pikanto qua kreskas sur ula planti. - DE.}
\l{}
\l{\sou{dorso}.\cel{1}\sou{(anat}.)Facio dopa di la homo-korpo ; facio supera di la}
\l{\cel{3}animalo-korpo. - EFIS.}
\l{}
\l{doselo. Kronatro qua superesas altaro, trono, e quan onu portas}
\l{\cel{3}en la procesioni super la eukaristio. - S.}
\l{}
\l{dotar.\cel{2}(\sou{trans., per, ye)}.\cel{2}Provizar (homo) per la havajo quan lu}
\l{\cel{3}adportas mariajo a sua spozo. -\cel{2}Provizar per ula qualesi.}
\l{\cel{3}- EFIS.}
\l{}
\l{dozo. I. Quanto de medikamento, quan la malado devas absorbar en}
\l{\cel{3}un foyo ; quanto quan onu devas inkluzar, en medikamento kom-}
\l{\cel{3}pozaja, de singla ek la substanci ye qui lu konsistas. - III.}
\l{\cel{3}Quanto de to quo formacas kompozajo. - DEFIRS.}
}
